\section{Abelian Categories}

\subsection{Zero morphisms, zero objects}

\bdefn

    A category $\C$ is said to have {\emphcolor zero morphisms} if for every $X,Y\in\C$ there exists a morphism $0_{XY}\colon X\to Y$ such that
    for all morphisms $f\colon X\to Y,g\colon Y\to Z$,
    $$ 0_{YZ}\circ f = g\circ0_{XY} = 0_{XZ} $$

\edefn

Notice that the zero morphism $0_{XY}$ is unique per $X,Y$.
Generally we omit the subscripts and just write e.g.\ $0\circ f=g\circ 0=0$.

\bdefn

    A {\emphcolor zero object} in a category is an object which is both terminal and initial.

\edefn

Clearly if a category has a zero object, all terminal objects are initial and vice versa.

\bprop

    A category with a zero object has zero morphisms.

\eprop

\bproof

    Let $0\in\C$ be the zero object, then for every $X,Y\in\C$ there is a unique map $0_{XY}\colon X\to Y$ which factors uniquely through $0$:
    $$ 0_{XY} = X\xto{\ !\ }0\xto{\ !\ }Y $$
    These are clearly zero morphisms.
    \qed

\eproof

\bexam

    The category of groups has a zero object: the trivial group.

\eexam

Consider a diagram $K\colon I\to\C$ which has both a limit and a colimit.
Recall
$$ \hom(\colim_I K,\lim_I K) \cong \lim_{i\in I^\op}\lim_{j\in I}\hom(Ki,Kj) \cong \lim_{(i,j)\in I^\op\times I}\hom(Ki,Kj) $$
Now we know that for a general functor $F\colon J\to\Set$, we have
$$ \lim F \cong \set{(a_j\in Fj)_j}[\forall\phi\colon i\to j.\,a_j=F(\phi)a_i] $$
In our case we have that for $(i,i'),(j,j')\in I^\op\times I$, a morphism $(i,j)\to(i',j')$ is a pair of morphisms $\phi\colon i'\to i$ and $\psi\colon j\to j'$.
The value of this functor on $(\phi,\psi)$ is the map $\hom(Ki,Kj)\to\hom(Ki',Kj')$ which maps $f\colon Ki\to Ki'$ to $K\psi\circ f\circ K\phi$.
Thus we have
$$ \hom(\colim K,\lim K) \cong \set{(f_{ij}\colon Ki\to Kj)_{ij}}[\forall\phi\colon i'\to i,\psi\colon j\to j'.\,f_{i'j'}=K\psi\circ f_{ij}\circ K\phi] $$

If $\C$ has zero morphisms it makes sense to consider the collection $(\delta_{ij}\colon Ki\to Kj)_{ij}$ defined by $\delta_{ii}={\rm id}_{Ki}$ and $\delta_{ij}=0$ when $i\neq j$.
Does this correspond to a morphism $\colim k\to\lim K$?
Let us consider when $i=j,i'=j'$, then we must have
$$ K\psi\circ K\phi = 1 \implies K(\psi\circ\phi) = 1 $$
for every $\phi\colon i'\to i,\phi\colon i\to i'$.
So $K$ must map endomorphisms to the identity (i.e.\ for $\phi\colon i\to i$, $K\phi=1$).
And when $i=j,i'\neq j'$ we must have $K(\psi\circ\phi)=0$.
So $K$ must map non-endomorphisms to $0$.

Such a $K$ can be viewed as a functor from the discrete category of $I$ (one with the same objects, but no non-identities), and as such $\delta$ is really on a morphism
when the diagram is from a discrete category, i.e.\ the limit and colimit are products and coproducts respectively.
But in such a case clearly $\delta$ is a morphism from $\coprod Ki$ to $\prod Ki$.
As we saw, morphisms $\coprod Ki\to\prod Ki$ can be viewed as (potentially infinite) matrices $(f_{ij}\colon Ki\to Kj)$, and $\delta$ can be viewed as the identity morphism.

\bdefn

    The {\emphcolor biproduct} of a collection of objects $\set{X_i}_{i\in I}$ is the product $\prod_iX_i$ such that $\coprod_iX_i$ also exists and the map $\delta\colon\coprod_iX_i\to\prod_iX_i$
    is an isomorphism.
    A biproduct is denoted $\bigoplus_iX_i$.

\edefn

The following is clear:

\bprop

    A collection $\set{X_i}_{i\in I}$ has a biproducts iff there is an object $\bigoplus_iX_i$ along with maps $p_j\colon\bigoplus_iX_i\to X_j$ and $i_j\colon X_j\to\bigoplus_iX_i$
    such that
    \benum
        \item $p_k\circ i_k=1_{X_k}$,
        \item $p_\ell\circ i_k=0$ for $\ell\neq k$,
        \item $\bigoplus_iX_i$ equipped with the $p_k$s is a product of $\set{X_i}_i$,
        \item $\bigoplus_iX_i$ equipped with the $i_k$s is a coproduct of $\set{X_i}_i$.
    \eenum

\eprop

\bdefn

    A category is {\emphcolor semiadditive} if it has zero morphisms and all binary biproducts.

\edefn

Suppose $\C$ is semiadditive, then for every $X\in\C$ we have canonical morphisms $\Delta\colon X\to X\oplus X$ and $\nabla\colon X\oplus X\to X$ given by the universal properties
of products and coproducts respectively.
That is, $p_k\circ\Delta=1_X$ for $k=1,2$ and $\nabla\circ i_k=1_X$ for $k=1,2$.

\bprop

    Let $\C$ be semiadditive, then every hom-set forms a commutative monoid such that composition is biadditive.

\eprop

\bproof

    Consider the operation defined by
    $$ \hom(X,Y)\times\hom(X,Y) \cong \hom(X,Y\oplus Y) \xto{\ \nabla_Y\ } \hom(X,Y) $$
    that is, we map $(f,g)$ to $\nabla_Y\circ(f,g)$, i.e.\ $f+g=\nabla(f,g)$.
    We claim that this gives hom-sets the structure of a commutative monoid.
    Firstly, this is commutative since the isomorphism is symmetric.
    Now consider
    $$ \hom(X,Y)^3 \cong \hom(X,Y\oplus Y\oplus Y) \xto{\ \nabla\ }\hom(X,Y) $$
    Since $Y\oplus(Y\oplus Y)\cong(Y\oplus Y)\oplus Y$, this is just the same as both $\bul+(\bul+\bul)$ and $(\bul+\bul)+\bul$, i.e.\ addition is associative.
    And we have $\nabla\circ(f,0)=\nabla\circ i_1\circ f=f$, so $0$ is the additive identity; i.e.\ the hom-sets indeed form commutative monoids.

    Now we need to show that $(f+g)\circ h=f\circ h+g\circ h$.
    This is because $h\circ (f+g)=h\circ \nabla\circ(f,g)$, and $h\circ\nabla\circ i_k=h$ meaning $h\circ\nabla=\nabla(h,h)$.
    Thus $h\circ(f+g)=\nabla(h,h)(f,g)=\nabla(hf,hg)=hf+hg$ as desired.
    \qed

\eproof

\subsection{Preadditive categories}

\bdefn

    A category $\C$ is {\emphcolor preadditive} if every hom-set is equipped with an operation $+$ turning it into an Abelian group, such that composition is bilinear.

\edefn

A preadditive category necessarily has zero morphisms: let $0_{XY}$ be the additive identity in $\hom(X,Y)$.
This is a zero morphism since $f\circ0=f\circ(0+0)=f\circ0+f\circ0$ so $f\circ0=0$, and similar for $0\circ f$.

\bprop

    All finite products and coproducts in a preadditive category are necessarily biproducts.

\eprop

\bproof

    It is sufficient to show this for binary products.
    Suppose $X\times Y$ exists, then define $i_X\colon X\to X\times Y$ by $i_X=(1_X,0)$ and similarly $i_Y=(0,1_Y)$.
    These form a coproduct: suppose $f\colon X\to Z$ and $g\colon Y\to Z$ are morphisms then the unique morphism $\phi\colon X\times Y\to Z$
    making the diagram commute is $\phi=f\circ p_X+g\circ p_Y$.
    Indeed, notice
    $$ \phi\circ(1_X,0) = f\circ p_X(1_X,0) + g\circ p_Y(1_X,0) = f + g\circ0 = f $$
    And if $\phi$ makes this diagram commute (i.e.\ $\phi\circ(1,0)=f$ and $\phi\circ(0,1)=g$) then
    $$ fp_X + gp_Y = \phi((1,0)p_X + (0,1)p_Y) $$
    and $(1,0)p_X+(0,1)p_Y$ is the identity (the identity of $X\times Y$ is the unique morphism such that $p_X1=p_X$ and $p_Y1=p_Y$, since morphisms are determined by
    their projections), so we see that $\phi=fp_X+gp_Y$ is the unique morphism satisfying this.

    So we have that $X\times Y$ is a coproduct as well, and its inclusion and projection maps are related by $p_k\circ i_\ell=\delta_{k\ell}$, so it is a biproduct.
    \qed

\eproof

\bdefn

    An {\emphcolor additive category} is a preadditive category which has all finite biproducts.

\edefn

Equivalently, an additive category is a preadditive category with all finite products, or all finite coproducts, or binary products and a terminal object, etc.
But importantly an additive category has binary biproducts and a zero object.

Notice that an additive category is a semiadditive category, and as such each hom-set is a commutative monoid as discussed above.
The commutative monoid structure coincides with the preadditive one.

\subsection{(Co)kernels, monomorphisms, epimorphisms}

\bdefn

    Let $\C$ be a category with zero morphisms.
    The {\emphcolor kernel} of a morphism $f\colon X\to Y$ is the equalizer $\ker f=\eq(f,0_{XY})$, and its {\emphcolor cokernel} is the coequalizer $\coker f=\coeq(f,0_{XY})$.

\edefn

Kernels and cokernels need not exist.
Explicitly the kernel of $f\colon X\to Y$ is a pair $(K\in\C,k\colon K\to X)$ such that
\benum
    \item $f\circ k=0$
    \item For any other $h\colon Z\to X$ such that $f\circ h=0$, there is a unique $\phi\colon Z\to K$ such that $h=k\circ\phi$.
\eenum
And the cokernel of $f\colon X\to Y$ is a pair $(Q\in\C,q\colon Y\to Q)$ such that
\benum
    \item $q\circ f=0$
    \item For any other $h\colon Y\to Z$ such that $h\circ f=0$, there is a unique $\psi\colon Q\to Z$ such that $h=\psi\circ q$.
\eenum

\bexam

    Consider $\grp$, which has a zero object and is complete and cocomplete.
    Thus $\grp$ has kernels and cokernels.
    The kernel of $f\colon G\to H$ is the usual group-theoretic kernel equipped with inclusion.
    The cokernel of $f$ is the quotient of $H$ by the normal closure of ${\rm im}f$, equipped with projection.

\eexam

\bprop

    Let $\C$ be preadditive with a zero object, then for every $f,g\colon X\to Y$, $\eq(f,g)=\ker(f-g)$ and $\coeq(f,g)=\coker(f-g)$.

\eprop

\bproof

    Notice that $h\circ(f-g)=0$ iff $h\circ f=h\circ g$.
    So cones over $f-g,0$ are precisely cones over $f,g$.
    \qed

\eproof

\bdefn

    A {\emphcolor monomorphism} is a morphism $f\colon X\to Y$ which is left-cancellative: $f\circ g=f\circ h$ implies $g=h$.
    An {\emphcolor epimorphism} is a morphism $f\colon X\to Y$ which is right-cancellative $g\circ f=h\circ f$ implies $g=h$.

\edefn

\bexam

    In $\Set$, monomorphisms are precisely injections, and epimorphisms are precisely surjections.

\eexam

\bprop

    Let $\C$ be a preadditive category with a zero object, and consider a morphism $f\colon X\to Y$.
    Let $(K,k\colon K\to X)$ be the kernel of $f$ and $(Q,q\colon Y\to Q)$ its cokernel.
    \benum
        \item $k$ is a monomorphism and $q$ is an epimorphism.
        \item If $\phi\colon Y\to Z$ is a monomorphism then $(K,k)$ is the kernel of $\phi\circ f$.
        Similarly if $\psi\colon Z\to X$ is an epimorphism then $(Q,q)$ is the cokernel of $f\circ\psi$.
        \item $f$ is a monomorphism iff $K$ is the zero object.
        And $f$ is an epimorphism iff $Q$ is the zero object.
    \eenum

\eprop

\bproof

    \benum
        \item First suppose $k\circ h=0$.
        Clearly $f\circ k\circ h=0$, and so by the universal property there is a unique morphism such that $k\circ h=k\circ g$.
        But clearly $g=0$ is such a morphism, and since it is unique this means that $h=0$.
        In general if $k\circ h=k\circ g$ then $k\circ(h-g)=0$ and so $h-g=0$ meaning $h=g$ as desired.

        \item Clearly $\phi\circ f\circ k=0$, so now we must just show it is universal.
        Suppose $\phi\circ f\circ h=0$, then $f\circ h=0$ by monomorphism, and so by the universal property of $\ker f$, $h$ factors through $k$.

        \item If $f$ is a monomorphism then $f\circ k=0$ implies $k=0$.
        This means that $K$ must be zero.
        Conversely, if $K=0$ (and so $k=0$), if $f\circ g=f\circ h$ then $f\circ(g-h)=0$ and so $g-h$ must factor through $0$ and as such $g-h=0$.
        So $g=h$ and $f$ is a monomorphism.
        \qed
    \eenum

\eproof

\bdefn

    A monomorphism is {\emphcolor normal} if it is the kernel of some morphism.
    Dually an epimorphism is {\emphcolor conormal} if it is the cokernel of some morphism.

\edefn

\bdefn

    A category $\C$ is {\emphcolor Abelian} if
    \benum
        \item it is preadditive with a zero object,
        \item has all binary biproducts (equivalently all finite biproducts),
        \item all kernels and cokernels exist,
        \item all monomorphisms are normal and all epimorphisms are conormal.
    \eenum

\edefn

An Abelian category has finite biproducts, equalizers, and coequalizers.
As such Abelian categories are necessarily finitely complete and finitely cocomplete.

\bexam

    The category $\ab$ is Abelian.
    Clearly it is preadditive, has a zero object, has binary biproducts, and kernels and cokernels exist ($\coker f=H/{\rm im}(f)$).
    Now we must show that monomorphisms are normal and epimorphisms are conormal.

    Let $f\colon A\to B$ be a monomorphism, then consider the projection $p\colon B\to B/{\rm im}(f)$ (the cokernel of $f$).
    Then $f$ is the kernel of $p$: clearly $p\circ f=0$ and if $p\circ h=0$ then we have ${\rm im}(h)\subseteq\ker p={\rm im}(f)$.
    $f$ is injective and so it has an inverse from its image $f^{-1}\colon{\rm im}(f)\to B$, and we have $f\circ f^{-1}\circ h=h$.

\eexam

In general one does have that monomorphisms satisfy $f=\ker\coker f$, and epimorphisms $f=\coker\ker f$.

Abelian categories are important as they allow one to do homology on them.
One can also discuss exact sequences in Abelian categories (to do so we must define what the image of a morphism is, which can be done in a general category, but in an Abelian category one
has that it is equal to $\ker\coker f$).

Another interesting result (but far too complicated to prove here) is the following:

\bthrm[title=Freyd-Mitchell]

    All Abelian categories are isomorphic to a full subcategory of the category of modules of some ring.

\ethrm

That is to say, for any Abelian category $\C$, there is a ring $R$ and a fully faithful embedding $\C\to{\rm Mod}_R$ (this embedding also preserves left and right exact sequences).

